%% Chapter 3 — Methodology
\label{ch:methodology}

This chapter presents the methodology behind the EMS Portal. The first section
analyzes the system requirements and introduces the three actors. The second
formalizes the multi-criteria matching algorithm, the project's core contribution.
The final section describes the engineering practices and principles used to build
the platform.

\section{Requirement Analysis}
The requirements are built around three actors: the \textbf{Client}, the
\textbf{Service Provider}, and the \textbf{Administrator}. They are split into
functional and non-functional needs.

\subsection{Functional Requirements}
The functional scope is delivered step by step across six phases. The
requirements are:

\begin{enumerate}
	\item \textbf{Authentication and identity.} Registration, login, logout, and
	      session refresh. Each user has an email, a name, a hashed password, a
	      role, and an account status.
	\item \textbf{Authorization.} Role-based access control for client, service
	      provider, and administrator, enforced at the route, permission, and
	      ownership levels.
	\item \textbf{Provider profiles and skills.} A service provider keeps a
	      discoverable profile (bio, location, hourly-rate range, and
	      availability) and attaches skills from a shared catalogue. The profile
	      has a clear publication lifecycle.
	\item \textbf{Service requests.} A client posts a request with the work,
	      budget, target location, and deadline. The request follows a status
	      state machine.
	\item \textbf{Multi-criteria matching.} The system ranks active, available
	      providers for a request using a configurable weighted scoring model, and
	      records the result for audit.
	\item \textbf{Engagements and ratings.} An accepted match becomes an
	      engagement. When it completes, both sides can rate each other, and
	      ratings feed back into a provider's reputation.
	\item \textbf{Administration.} An administrator manages users, tunes matching
	      weights, and reviews how well the matching policy works.
\end{enumerate}

\subsection{Non-Functional Requirements}
The platform must be \textbf{secure} (signed tokens, hashed secrets, enforced
authorization, and protection against common web vulnerabilities). It must be
\textbf{maintainable} (a modular, layered architecture with automated tests and
continuous integration). It must be \textbf{performant} (low request latency,
measured under load). And it must be \textbf{scalable} (a stateless API and an
indexed schema that supports the matching hot path).

\subsection{Use Case Description}
The three actors interact with the system in distinct ways.

\textbf{Client Workflow:} A Client registers and signs in, then browses
published provider profiles and their skills. When ready, the client posts a
service request describing the work, budget, location, and deadline. The system
returns a ranked list of matched providers. The client reviews the rankings and
their reasoning, starts an engagement with a chosen provider, and after the work
completes, rates the provider.

\textbf{Service Provider Workflow:} A Service Provider signs in and creates a
profile, specifying their bio, hourly-rate range, and availability. They then
manage their skills, selecting from the shared catalogue to showcase their
expertise. Once the profile is complete and publishable, they approve it for
discovery. Thereafter, they view incoming service requests, accept or decline
engagements, and receive ratings from clients they have worked with.

\textbf{Administrator Workflow:} An Administrator manages the user base (viewing,
suspending, or removing accounts), curates and updates the shared skill
catalogue, and tunes the matching algorithm. Specifically, the administrator can
set the weights for the five matching criteria, compare the results of different
weighting policies on historical requests, and audit how well the matching system
works.

Figure~\ref{fig:usecase} sums up these interactions.

\subsection{Use Case Diagram}
\begin{figure}[ht]
	\centering
	\begin{tikzpicture}[
		actor/.style={font=\footnotesize, align=center},
		usecase/.style={draw, ellipse, align=center, font=\scriptsize,
			minimum width=2.7cm, minimum height=0.85cm, fill=blue!5},
		boundary/.style={draw, rounded corners, inner sep=12pt},
		line/.style={-}
	]
		% Use cases (system boundary)
		\node[usecase] (login)   {Register /\\ Authenticate};
		\node[usecase] (profile) [below=0.35cm of login]   {Manage profile\\ \& skills};
		\node[usecase] (request) [below=0.35cm of profile] {Post service\\ request};
		\node[usecase] (match)   [below=0.35cm of request] {View ranked\\ matches};
		\node[usecase] (engage)  [below=0.35cm of match]   {Manage\\ engagement};
		\node[usecase] (rate)    [below=0.35cm of engage]  {Rate\\ counterpart};
		\node[usecase] (admin)   [below=0.35cm of rate]    {Administer \&\\ tune weights};

		% System boundary box around the use cases
		\begin{scope}[on background layer]
			\node[boundary, fit=(login)(profile)(request)(match)(engage)(rate)(admin),
				label=above:{\footnotesize\textbf{EMS Portal}}] (sys) {};
		\end{scope}

		% Actors
		\node[actor] (client) [left=2.2cm of match] {\textbf{Client}};
		\node[actor] (provider) [right=2.2cm of profile] {\textbf{Service}\\\textbf{Provider}};
		\node[actor] (adminA) [right=2.2cm of admin] {\textbf{Admin}};

		% Client associations
		\draw[line] (client) -- (login.west);
		\draw[line] (client) -- (request.west);
		\draw[line] (client) -- (match.west);
		\draw[line] (client) -- (engage.west);
		\draw[line] (client) -- (rate.west);

		% Provider associations
		\draw[line] (provider) -- (login.east);
		\draw[line] (provider) -- (profile.east);
		\draw[line] (provider) -- (engage.east);
		\draw[line] (provider) -- (rate.east);

		% Admin associations
		\draw[line] (adminA) -- (admin.east);
		\draw[line] (adminA) -- (login.east);
	\end{tikzpicture}
	\caption{High-level use case diagram for the three EMS Portal actors.}
	\label{fig:usecase}
\end{figure}

\subsection{System Overview}
The EMS Portal is a client-server web application with three main tiers. The
browser runs a single-page application (SPA) that handles user interaction and
renders the interface. The SPA communicates with a stateless REST API that
contains all business logic; the authorization layer sits within the API and
grants different capabilities to clients, providers, and administrators based on
their role. The API reads from and writes to a relational database, which serves
as the system of record for all user, profile, request, and engagement data.

The matching engine is a core component within the API. When a client posts a
service request, the engine reads the request details and the set of active,
available providers from the database. It scores each provider against the request
using the five criteria (skill, availability, cost, location, rating), weights
them according to a configurable matching policy, and produces a ranked list with
full scoring transparency. Both the input request and the per-provider scores are
logged to an audit table for transparency and analysis.

The actual architecture diagrams and implementation details are presented in
Chapter~\ref{ch:prototyping}.

\section{Algorithms}
The defining feature of the EMS Portal is its \textbf{multi-criteria matching
algorithm}. It ranks providers for a given service request. This section
formalizes the scoring model. The implementation is planned for Phase~4 (see the
work plan in Chapter~\ref{ch:conclusion}).

\subsection{Formalizing}
Let $R$ be a service request and $P$ the set of candidate providers that are
active and available. Each provider $p \in P$ is scored against $R$ on five
criteria, indexed by $i \in \{\text{skill}, \text{avail}, \text{cost},
\text{loc}, \text{rating}\}$. Each criterion gives a normalized score
$c_i(p, R) \in [0, 1]$, and a matching configuration supplies non-negative
weights $w_i$ with $\sum_i w_i = 1$. The total score of a provider is the
weighted sum

\begin{equation}
	\text{Score}(p, R) = \sum_{i} w_i \cdot c_i(p, R),
	\label{eq:score}
\end{equation}

and providers are ranked by descending $\text{Score}$. The five criterion scores
are defined as follows.

\begin{itemize}
	\item \textbf{Skill compatibility} ($c_\text{skill}$): the overlap between the
	      skills required by $R$ and the skills held by $p$, measured as the
	      fraction of required skills the provider has.
	\item \textbf{Availability} ($c_\text{avail}$): a normalized measure based on
	      the provider's availability flag and current engagement load.
	\item \textbf{Cost} ($c_\text{cost}$): how much the request's budget range
	      overlaps the provider's hourly-rate range.
	\item \textbf{Location proximity} ($c_\text{loc}$): a value that falls off
	      with the great-circle distance between the request's target location and
	      the provider, computed with the haversine formula in
	      Equation~\ref{eq:haversine}.
	\item \textbf{Rating} ($c_\text{rating}$): the provider's mean past rating,
	      normalized to $[0, 1]$.
\end{itemize}

The great-circle distance $d$ between two points given by latitude/longitude
$(\varphi_1, \lambda_1)$ and $(\varphi_2, \lambda_2)$ on a sphere of radius $r$
is

\begin{equation}
	d = 2r \cdot \arcsin\!\sqrt{\sin^2\!\Big(\tfrac{\varphi_2 - \varphi_1}{2}\Big)
	+ \cos\varphi_1 \cos\varphi_2 \sin^2\!\Big(\tfrac{\lambda_2 - \lambda_1}{2}\Big)}.
	\label{eq:haversine}
\end{equation}

The weights $w_i$ are not hard-coded. They are stored as named \textbf{matching
configurations}. Two configurations are seeded for comparison: \texttt{basic-v1},
which gives every criterion equal weight, and \texttt{intelligent-v1}, which uses
tuned weights. Running both on the same requests lets the intelligent policy be
measured directly against the simple baseline. Algorithm~\ref{alg:match} gives
the ranking procedure.

\begin{algorithm}[H]
	\small
	\caption{$(\text{Ranked}) \gets \texttt{MatchProviders}(R, W)$}
	\label{alg:match}
	\begin{algorithmic}[1]
		\Require Request $R$; matching configuration $W = \{w_i\}$ with $\sum_i w_i = 1$.
		\State $\text{Results} \gets \emptyset$
		\State $P \gets \texttt{ActiveAvailableProviders}()$
		\For{$p \in P$}
			\State $c_\text{skill} \gets \texttt{SkillOverlap}(R, p)$
			\State $c_\text{avail} \gets \texttt{Availability}(p)$
			\State $c_\text{cost} \gets \texttt{CostOverlap}(R, p)$
			\State $c_\text{loc} \gets \texttt{ProximityScore}(R, p)$ \Comment{via Eq.~\ref{eq:haversine}}
			\State $c_\text{rating} \gets \texttt{NormalizedRating}(p)$
			\State $s \gets \sum_{i} w_i \cdot c_i$ \Comment{Eq.~\ref{eq:score}}
			\State $\text{Results} \gets \text{Results} \cup \{(p, s, \{c_i\})\}$
		\EndFor
		\State \Return $\texttt{SortDescendingByScore}(\text{Results})$
	\end{algorithmic}
\end{algorithm}

Each row of the result, the per-criterion scores and the total, is saved to an
audit table. This lets a ranking be explained, and lets different configurations
be compared afterward.

\section{Development Methodology}
\label{sec:dev-method}
The platform is built with a clear, repeatable methodology that enforces
consistency across boundaries, enables reliable refactoring, and ensures both
the codebase and database schema can be confidently versioned and reproduced.

\begin{itemize}
	\item \textbf{Single source of truth for types.} Cross-boundary data shapes
	      are defined once as Zod schemas in a shared package and used by both
	      ends. Backend entities \texttt{implement} these types, so a schema change
	      forces a compile error wherever the shape is used.
	\item \textbf{Layered modules with the Data Mapper pattern.} Each backend
	      feature is a module layered as controller $\rightarrow$ service
	      $\rightarrow$ repository $\rightarrow$ entity. Queries live in injectable
	      repositories, not in the entities.
	\item \textbf{Migrations-first database.} Schema synchronization is off. Every
	      schema change is a reviewed, versioned migration, applied on purpose.
	\item \textbf{Testing pyramid.} A large set of fast, isolated unit tests is
	      backed by a smaller set of integration tests that run the API against a
	      real database, with enforced coverage thresholds.
	\item \textbf{Continuous integration.} Every change passes style, build, unit,
	      security, integration, and deployment-readiness checks before it can
	      reach the main branch.
\end{itemize}

These practices are described in concrete, implemented terms in
Chapter~\ref{ch:prototyping}.
